\documentclass[11pt]{article}
\usepackage{amsmath,amssymb,amsthm,geometry,hyperref}
\geometry{margin=1in}
\newtheorem{theorem}{Theorem}
\newtheorem{lemma}[theorem]{Lemma}
\newtheorem{proposition}[theorem]{Proposition}
\newcommand{\ddc}{dd^c}
\newcommand{\PP}{\mathbb P}
\newcommand{\OO}{\mathcal O}
\newcommand{\II}{\mathcal I}
\newcommand{\ord}{\operatorname{ord}}
\title{An explicit realization criterion for arbitrary formal multiplier-ideal data on the projective line}
\author{Working mathematical draft: continuation 308}
\date{8 October 2026}
\begin{document}
\maketitle
\begin{abstract}
For an arbitrary family of nonzero coherent ideal subsheaves of the structure sheaf of the complex projective line, we give a necessary and sufficient criterion for realization as all multiplier ideals of one quasi-subharmonic potential with prescribed Fubini--Study curvature budget. The criterion uses only the orders of the supplied ideals, a countable supremum and a nonnegative series; it contains no unknown potential or multiplier-ideal operation. An explicit series of logarithmic chordal kernels realizes every admissible family. The local integrability proof includes the critical integer exponents and uses no strong-openness theorem. This is a complete proposed theorem on the projective line, pending independent review. It is a scoped result: the arbitrary-dimensional original P14.8 remains open, and candidate status is false.
\end{abstract}

\section{Conventions and the data-only criterion}
Let $X=\PP^1(\mathbb C)$ and use the convention
\[
 dd^c=\frac{i}{2\pi}\partial\bar\partial,
 \qquad dd^c\log|z|^2=\delta_0.
\]
Write $\omega=dd^c\log(1+|z|^2)$ for the Fubini--Study form, with $\int_X\omega=1$, and fix a real number $V\geq0$. A potential means an upper semicontinuous $u\in L^1(X)$, not identically $-\infty$, for which $V\omega+dd^cu$ is a positive measure. The multiplier-ideal convention throughout is
\[
 \II(cu)_p=\{f\in\OO_{X,p}: |f|^2e^{-cu}\text{ is locally integrable near }p\},
 \qquad c>0.
\]
It uses $e^{-cu}$, rather than $e^{-2cu}$.

The input is an arbitrary family $(A_c)_{c>0}$ of nonzero coherent ideal subsheaves $A_c\subseteq\OO_X$. It is not assumed realizable, monotone or continuous in the parameter. On a smooth compact curve each such ideal is $\OO_X(-D_c)$ for a finite effective integral divisor. Equivalently, its stalk at $p$ is generated by a local parameter to the nonnegative integer order
\[
 k_c(p)=\ord_p A_c.
\]
Set, entirely from the integer-indexed supplied data,
\[
 S_0=\bigcup_{m\geq1}\operatorname{Supp}D_m,
 \qquad b_p=\sup_{m\geq1}\frac{k_m(p)}m \quad(p\in X).
\]
The set $S_0$ is countable and $b_p=0$ off $S_0$. Infinite values of $b_p$ are allowed at this preliminary stage. Sums of nonnegative values mean the supremum of all finite subsums.

\begin{theorem}\label{main}
The family $(A_c)_{c>0}$ is realized by one $V\omega$-psh potential $u$, in the sense that $\II(cu)=A_c$ for every real $c>0$, if and only if the following two conditions hold:
\begin{align}
 &b_p<\infty\text{ for all }p,
 \qquad B:=\sum_{p\in S_0}b_p\leq V; \label{budget}\\
 &k_c(p)=\lfloor c b_p\rfloor
 \quad\text{for every }p\in X\text{ and every real }c>0. \label{floor}
\end{align}
In particular this is an explicit test on arbitrary formal ideal data, not a restatement involving an unknown realizing function. Whenever these conditions hold a realizing potential is
\[
 u(Z)=\sum_{p\in S_0}b_p\log
 \frac{|\det(Z,P_p)|^2}{\|Z\|^2\|P_p\|^2},
\]
where $Z$ and $P_p$ are nonzero homogeneous representatives. The nonpositive series is understood as its decreasing limit. It has $\int_X(-u)\omega=B$, and its curvature is
\[
 V\omega+dd^cu=\sum_{p\in S_0}b_p\delta_p+(V-B)\omega.
\]
A finite additive constant can normalize its supremum to zero without changing the ideals.
\end{theorem}

The criterion has infinitely many tests; no finite decision algorithm or arbitrary-dimensional analogue is asserted. Even the formal orders at noninteger indices must pass \eqref{floor}. The budget alone is insufficient.

\section{The local integrability calculation}
We use elementary distributional potential theory in one complex variable: the logarithmic fundamental solution, Fubini's theorem for its locally integrable kernel, and Weyl's lemma that a harmonic distribution is a smooth harmonic function. Equivalently these imply the usual local Riesz representation. The following argument states precisely the use of these ordinary analytic facts and supplies the needed integrability estimates.

\begin{lemma}\label{local}
Let $w$ be subharmonic and not identically $-\infty$ near $0$. Let $a$ be the atom at $0$ of its positive measure $dd^cw$. For a nonzero holomorphic germ $f$ of order $k$ and a real $c>0$,
\[
 |f|^2e^{-cw}\in L^1_{\mathrm{loc}}(0)
 \quad\Longleftrightarrow\quad k+1>ca.
\]
Consequently $\II(cw)_0=(z^{\lfloor ca\rfloor})$.
\end{lemma}
\begin{proof}
Remove the atom distributionally: $w=a\log|z|^2+v$, where $v$ has a subharmonic representative with atom-free curvature at $0$. The equality is sufficient almost everywhere for the integrability question. This removal follows either from local Riesz representation or directly by subtracting the fundamental solution, retaining the positive residual curvature and its upper semicontinuous representative. On a smaller disk $v$ is bounded above. Write $f=z^kg$ with a holomorphic unit $g$, bounded above and away from zero on that disk. If $k+1\leq ca$, an upper bound for $v$ gives a positive lower bound by a constant times $|z|^{2k-2ca}$. Its radial area integral diverges, including the equality case. This proves necessity without omitting a critical exponent.

For sufficiency put $q=\max(ca-k,0)<1$. Choose a small disk $D_r$ whose residual curvature mass $m$ satisfies $cm<1-q$; this is possible because there is no atom at $0$. On a smaller concentric disk the local representation is
\[
 v(z)=\int_{D_r}\log|z-\zeta|^2\,d\mu(\zeta)+h(z),
 \qquad \mu(D_r)=m,
\]
with $h$ harmonic and bounded there. Indeed the logarithmic integral is locally $L^1$ by Fubini, has distributional curvature $\mu$, and subtracting it from $v$ leaves a harmonic distribution. The representation need only hold almost everywhere.

If $m=0$, boundedness of $h$ reduces integrability to $q<1$. If $m>0$, Jensen's inequality for the probability measure $\mu/m$, first with truncated logarithms and then by monotone convergence, gives
\[
 \exp\left(-c\int\log|z-\zeta|^2\,d\mu(\zeta)\right)
 \leq\frac1m\int |z-\zeta|^{-2cm}\,d\mu(\zeta).
\]
Choose conjugate finite H\"older exponents $\alpha,\beta>1$ such that $q\alpha<1$ and $cm\beta<1$. They exist precisely because $q+cm<1$; when $q=0$, first choose $1<\beta<1/(cm)$. H\"older's inequality gives a uniform bound, for $\zeta\in D_r$, on
\[
 \int_{D_{r/2}} |z|^{-2q}|z-\zeta|^{-2cm}\,dA(z).
\]
The first factor has finite norm because $q\alpha<1$; the second has uniformly finite norm because its translated integration region lies in a disk of radius $2r$ and $cm\beta<1$. Integrate this bound against $\mu/m$, using Tonelli. Bounded harmonic factors and the unit $g$ do not alter local integrability. This proves sufficiency. The smallest nonnegative integer $k$ with $k+1>ca$ is $\lfloor ca\rfloor$, proving the ideal assertion.
\end{proof}

The same statement applies to quasi-subharmonic potentials after adding a smooth local curvature potential. A smooth addition multiplies the integrand by a bounded positive factor. Thus the relevant atom is that of the positive curvature measure, without a change of the exponent convention.

\section{Necessity of both formal-data conditions}
Suppose $u$ realizes the family and put $\nu=V\omega+dd^cu$. Distributional Stokes, namely testing $dd^cu$ against the constant function one, gives $\nu(X)=V$. Write $a_p=\nu(\{p\})$. These atoms are nonnegative, countably supported, and their sum is at most $V$.

Near any point add $V\log(1+|z|^2)$ in a Fubini--Study coordinate to $u$. The resulting function is subharmonic with curvature $\nu$. Lemma~\ref{local} gives
\[
 k_c(p)=\lfloor c a_p\rfloor \qquad(c>0).
\]
For any finite $a\geq0$,
\[
 \sup_{m\geq1}\frac{\lfloor ma\rfloor}{m}=a,
\]
since the ratios are at most $a$ and are greater than $a-1/m$. Hence the $b_p$ computed from the supplied integer ideals equal $a_p$. Every point with $a_p>0$ belongs to $S_0$, by choosing an integer $m>1/a_p$. The budget \eqref{budget} and the whole real-parameter identity \eqref{floor} follow. No smoothness or finiteness of the pole set of $u$ was assumed.

\section{Explicit sufficiency and the countable-pole construction}
For $p\in X$ define its chordal logarithm
\[
 \phi_p(Z)=\log\frac{|\det(Z,P_p)|^2}{\|Z\|^2\|P_p\|^2}\leq0.
\]
Cauchy--Schwarz gives the inequality and makes the expression independent of representatives. The local logarithmic fundamental solution gives
\[
 dd^c\phi_p=\delta_p-\omega.
\]
Its $L^1$ size is uniformly and exactly
\[
 \int_X(-\phi_p)\omega=1.
\]
To verify the last identity, unitary invariance sends $p$ to infinity; in the affine chart $\phi_\infty=-\log(1+|z|^2)$ and the radial Fubini--Study density is $2r(1+r^2)^{-2}dr$. The integral is $\int_0^1-\log t\,dt=1$ after $t=(1+r^2)^{-1}$.

Assume \eqref{budget}--\eqref{floor}. Enumerate the positive $b_p$, allowing a finite or empty list, and take finite partial sums $u_N=\sum_{j\leq N}b_{p_j}\phi_{p_j}$. They decrease, and
\[
 V\omega+dd^cu_N=\sum_{j\leq N}b_{p_j}\delta_{p_j}
 +\left(V-\sum_{j\leq N}b_{p_j}\right)\omega\geq0.
\]
By Tonelli, the decreasing series has $\int(-u)\omega=B<\infty$, so it is not identically $-\infty$ and $u_N\to u$ in $L^1$. Its pointwise decreasing limit is upper semicontinuous. Adding the smooth local Fubini--Study potential, use the disk-area version of the submean inequalities: these pass to the decreasing limit, since local $L^1$ convergence controls the disk integrals. Together with upper semicontinuity this is the subharmonic representative, and $u$ is $V\omega$-psh. Alternatively its positive distributional curvature and the same decreasing representative establish this conclusion.

Distributional differentiation of the $L^1$ limit gives
\[
 V\omega+dd^cu=\sum_p b_p\delta_p+(V-B)\omega.
\]
The measure series converges in total variation because its nonnegative total mass is finite. The atom at any $p$ is exactly $b_p$, even at accumulation points of other poles; a countable sum of point masses assigns to a singleton precisely its own coefficient. Lemma~\ref{local} therefore gives $\II(cu)_p=(z^{\lfloor cb_p\rfloor})=(A_c)_p$ at every point and every $c>0$. Equality of stalks proves equality of ideal sheaves, completing Theorem~\ref{main}.

The divisors prescribed by the criterion are finite: if $\lfloor cb_p\rfloor>0$ then $b_p\geq1/c$, so there are at most $cB$ such distinct points and the total integral multiplicity is at most $cB$. This verifies coherence rather than assuming it for the constructed family. The empty family of positive coefficients gives $u=0$. If $V=0$, the budget forces every $b_p=0$ and every ideal to be $\OO_X$. Finally $u\leq0$ and $\int(-u)\omega=B$ imply $-B\leq\sup_Xu\leq0$; its finite supremum may be subtracted for normalization.

\section{Two formal obstructions and the remaining general problem}
The two conditions are distinct. If $V=1$ and $A_c=\OO_X(-\lfloor2c\rfloor p)$ at one point, then the exact linear floor pattern holds with $b_p=2$, but the curvature budget fails. Hence no $\omega$-psh potential realizes it.

Conversely let $V=2$ and prescribe at one point the orders $0$ for $0<c<1$, $1$ for $1\leq c<2$, and $3$ for $c\geq2$, with order zero elsewhere. This is a monotone, right-continuous family of coherent ideals. Integer data give $b_p=3/2$, within the budget, but the required order at $c=4$ would be six, not three. Thus even monotonicity, one-sided constancy and the mass budget do not suffice; the full parameter compatibility matters.

Realization does not recover the raw potential. For $V>0$, both zero and sufficiently small nonconstant smooth functions $v$ with $V\omega+dd^cv>0$ have all multiplier ideals equal to $\OO_X$. Such $v$ exist by bounding the curvature of any fixed nonconstant smooth function against $\omega$ and scaling it. Their curvature measures can differ. The theorem neither identifies a canonical normalized $\II$-model nor asserts descent of arbitrary representative invariants through ideal data.

The result treats arbitrary formal data, not merely already-realizable inputs, but only on $(\PP^1,V\omega)$. It does not settle arbitrary-dimensional compact K\"ahler spaces, big non-nef classes, all-semivaluation reconstruction, non-pluripolar mass descent or the remaining zero-mass interior range in original P14.8. No native active campaign output was used in this independent working draft, and none is fed this draft while active. It remains unreviewed; no original closure or new accepted-paper count is asserted. The protected P16.4 editor is unchanged.
\end{document}
